\section{Experiments and Results}
\label{sec:experiments_results}

The experiments were organised around the main empirical tension of this study: the cohort contained large-scale multimodal evidence, but report-level semantic supervision was sparse, centre-dependent, and therefore unsuitable for direct fully supervised report-guided learning. We therefore evaluated LCAD--RASA through a progressive sequence of analyses. We first quantified the supervision bottleneck, then assessed whether pseudo-reports contained modality-grounded information rather than label-derived templates. We next tested whether these structured reports provided measurable cross-modal alignment signals and whether such signals translated into improved prediction under report scarcity. Finally, we examined the robustness of the observed behaviour through ablation, perturbation, reference-stratified testing, and leave-one-centre-out evaluation.

\subsection{A five-centre cohort revealed a semantic-supervision bottleneck rather than an imaging-scale bottleneck}
\label{subsec:cohort_results}

The final analytic cohort comprised 1,897 multimodal cervical screening cases from five centres, with 137,591 evaluable OCT and colposcopy images (Table~\ref{tab:cohort_summary}, ``Cohort scale and case-level report supervision''; Figure~\ref{fig:centre_supervision}, ``Centre-level cohort scale and report-supervision imbalance''). This scale indicates that image availability was not the principal constraint in the present dataset. The more restrictive bottleneck was the uneven availability of clinician-authored report supervision. Real archived reports were available for 744 cases, whereas 1,153 cases had complete multimodal evidence but no archived report and were therefore eligible for structured pseudo-report completion.

The imbalance in report supervision was strongly centre-dependent. Enshi retained real reports for all 406 cases, Jingzhou retained reports for 334 of 406 cases, Xiangyang retained only 4 reports among 500 cases, and Wuhan and Shiyan had no centre-wide real-report archives. This distribution defines the learning problem addressed by LCAD--RASA. A report-aware model cannot assume that semantic supervision is uniformly observed across centres; it must instead learn from abundant image and clinical evidence while using incomplete report supervision in a controlled and auditable way. The next analysis therefore asked whether pseudo-reports could serve as structured semantic anchors, rather than merely converting diagnostic labels into formatted text.

\subsection{Structured pseudo-reports provided modality-grounded semantic anchors rather than label templates}
\label{subsec:pseudo_report_results}

To determine whether pseudo-reports carried modality-specific information, we compared structured pseudo-report sources against a label-template baseline. The label-template baseline served as a negative control because it could satisfy the required report schema without using image or clinical evidence. As expected, it was schema-complete but had no modality support, with a mean modality-support rate of 0.000, and exhibited substantial textual repetition, with a maximum duplicate fraction of 0.867 (Table~\ref{tab:pseudo_source_comparison}, ``Structured pseudo-report source comparison''; Figure~\ref{fig:pseudo_source_comparison}, ``Structured pseudo-report source comparison across template, rule-based, and local LLM sources''). This result demonstrates that schema validity alone is insufficient: a report-like output may be structurally complete while still failing to encode OCT, colposcopy, or clinical context.

In contrast, both the rule-based and local embedding-assisted LLM pseudo-reports preserved OCT, colposcopy, and clinical section support, with a mean modality-support rate of 1.000. The rule-based baseline achieved a label-consistency mean of 0.628, and the local LLM baseline achieved the same label-consistency mean while maintaining modality-grounded sections. These results suggest that structured pseudo-reports can provide semantically organised evidence for representation learning, rather than functioning as simple textual restatements of class labels.

We further examined whether pseudo-report generation was dependent on a single language-model provider. Template, rule-based, local embedding LLM, Qwen-Plus, GLM-4.7-Flash, Gemini-3.1-Pro, and GPT-5.5 were evaluated under a common structured-report schema (Table~\ref{tab:llm_provider_comparison}, ``LLM provider comparison for structured pseudo-report generation''; Figure~\ref{fig:api_provider_summary_heatmap}, ``LLM provider comparison for structured pseudo-report generation''; Figure~\ref{fig:api_quality_heatmap}, ``Stage 1 provider quality metrics and risk-oriented quality indicators''; Figure~\ref{fig:api_reliability}, ``Provider latency, modality support, and generation reliability''). All evaluated sources produced schema-valid structured outputs in their available cohorts. GPT-5.5 was evaluated on the extended 100-case cohort and generated 100 valid structured reports, with a schema-valid rate of 1.000, section completeness of 0.810, modality support of 0.810, contradiction rate of 0.050, hallucination rate of 0.000, mean latency of 11.39 seconds per case, and an estimated cost of US\$27.73 per 1,000 cases. Qwen-Plus, GLM-4.7-Flash, and Gemini-3.1-Pro were retained as pilot cohorts; their results are therefore interpreted as feasibility evidence, not as a definitive head-to-head ranking of providers.

Only de-identified structured evidence was submitted to external APIs. Raw images, image paths, patient names, internal patient identifiers, and hospital identifiers were not transmitted. These safeguards preserve the methodological distinction between pseudo-report augmentation as a controllable learning component and clinician-authored reporting as a clinical record. Having established that pseudo-reports can contain modality-grounded sections, we next examined whether those sections were aligned with the corresponding multimodal representations.

\subsection{Report sections exposed measurable cross-modal semantic alignment}
\label{subsec:alignment_results}

We evaluated report-guided alignment by matching each modality representation to its corresponding report section: OCT embeddings to \texttt{oct\_findings}, colposcopy embeddings to \texttt{colposcopy\_findings}, clinical embeddings to \texttt{clinical\_context}, and fused representations to \texttt{impression}. Full LCAD--RASA achieved a macro mean reciprocal rank (MRR) of 0.049 across the four sections. Removing section alignment reduced macro MRR to 0.022, while simple concatenation fusion remained similarly low at 0.020 (Table~\ref{tab:theme1_alignment_ablation}, ``Direct RASA alignment ablation using modality-section retrieval''; Figure~\ref{fig:alignment_retrieval_mrr}, ``Modality-section retrieval MRR for Theme 1 cross-modal semantic alignment'').

The highest retrieval-only macro MRR was observed for the real-report-only variant, at 0.056. This result is informative because it separates the quality of semantic targets from the coverage of semantic supervision. Real reports provide strong alignment targets when they are available, but their limited and centre-imbalanced distribution constrains their usefulness as the sole source of report supervision. LCAD--RASA therefore does not treat pseudo-reports as replacements for clinician-authored reports. Instead, it uses structured pseudo-reports to extend report-guided alignment to cases in which multimodal evidence exists but archived report supervision is missing.

The staged API alignment analysis provided a complementary view of provider-dependent alignment behaviour. Gemini-3.1-Pro achieved the highest available pilot macro MRR of 0.366, whereas GPT-5.5 achieved a macro MRR of 0.063 on the extended cohort (Figure~\ref{fig:api_alignment_mrr}, ``Stage 2 provider alignment analysis: macro MRR and section-specific MRR''; Figure~\ref{fig:api_section_mrr}, ``Provider-level section MRR across report sections''). Because these evaluations differed in cohort size and stage, the values should not be interpreted as a direct provider ranking. They nevertheless show that structured report sections provide explicit and measurable alignment targets across pseudo-report sources. The subsequent analysis therefore asked whether these alignment signals were most useful in the intended setting: limited availability of real reports.

\subsection{Pseudo-report augmentation was most beneficial when real reports were scarce}
\label{subsec:scarcity_results}

To test whether pseudo-report augmentation addressed the report-supervision bottleneck, we deliberately varied the availability of real reports and evaluated a lightweight downstream surrogate. The largest gain appeared in the lowest-supervision regime (Table~\ref{tab:scarcity_curve}, ``Report-supervision scarcity curve using a lightweight downstream surrogate''; Figure~\ref{fig:scarcity_curve}, ``Report-supervision scarcity curve comparing real-report-only and pseudo-report-augmented surrogates''). At 10\% real-report availability, the LCAD-augmented surrogate achieved AUROC 0.818 and F1 0.502, compared with AUROC 0.636 and F1 0.297 for the real-report-only surrogate. At full real-report availability, the augmented surrogate also remained higher, with AUROC 0.841 compared with 0.796 for the real-report-only surrogate.

This pattern indicates that pseudo-report augmentation is not merely adding redundant text when clinician reports are already available. Its main value emerges in the clinically realistic setting in which OCT, colposcopy, and structured clinical variables are present, but report-level supervision is incomplete. The staged API downstream surrogate was intentionally conservative. GPT-5.5 achieved AUROC 0.630 and F1 0.229 at 10\% real-report availability on the extended cohort, whereas Gemini-3.1-Pro achieved AUROC 0.725 and F1 0.379 in its pilot cohort (Figure~\ref{fig:api_scarcity_auc}, ``Stage 3 downstream scarcity surrogate for selected API-generated pseudo reports''). These analyses support the utility of structured pseudo-report augmentation under report scarcity, while remaining distinct from full end-to-end retraining evidence. We therefore proceeded to evaluate whether the full LCAD--RASA configuration improved held-out diagnostic performance under the locked retrospective protocol.

\subsection{Held-out testing favoured the full LCAD--RASA configuration in point estimates}
\label{subsec:diagnostic_results}

On the locked held-out test set ($n=288$), Full LCAD--RASA achieved AUROC 0.832 (95\% CI 0.757--0.897) and F1 0.611 (95\% CI 0.458--0.737) at the validation-selected threshold of 0.37 (Table~\ref{tab:main_comparison}, ``Held-out test performance at validation-selected thresholds''; Figure~\ref{fig:main_auc}, ``Held-out AUROC with bootstrap confidence intervals''; Figure~\ref{fig:main_metrics}, ``Held-out multi-metric profile across model variants''; Figure~\ref{fig:main_auc_f1}, ``AUROC--F1 trade-off across main model variants''). Pseudo-augmented LCAD training achieved AUROC 0.826 (95\% CI 0.747--0.894) and F1 0.545 (95\% CI 0.395--0.674). Real-report-only training achieved AUROC 0.725, whereas simple concatenation fusion achieved AUROC 0.472.

The ordering of these internal configurations is consistent with the preceding analyses. Generic concatenation was insufficient for multimodal fusion; real-report-only learning was limited by sparse and uneven report coverage; pseudo-report augmentation improved over real-report-only learning; and the full LCAD--RASA model produced the highest held-out point estimates among the report-aware internal variants. Corrected paired bootstrap rechecking supported clear AUROC differences between Full LCAD--RASA and weaker internal comparators such as simple concatenation fusion and real-report-only training, but not between Full LCAD--RASA and pseudo-augmented LCAD. The appropriate interpretation is therefore that Full LCAD--RASA produced favourable internal held-out performance under the locked retrospective protocol, rather than definitive evidence of clinical superiority.

\subsection{External same-split baseline comparison}
\label{subsec:external_baseline_results}

To determine whether the observed pattern depended on weak internal comparators, we added an external baseline block under the same train/validation/test split, validation-selected threshold rule, held-out test set, and bootstrap confidence-interval protocol. The block included clinical-only logistic regression, clinical-only gradient boosting, OCT-only and colposcopy-only frozen-embedding MLP classifiers, a late-fusion MLP, a cross-attention multimodal transformer, and a CLIP-style contrastive multimodal baseline without report-section supervision (Table~\ref{tab:external_baselines}, ``Same-split external baseline comparison on the locked held-out test set''; Figure~\ref{fig:external_baselines_auc}, ``Same-split external baseline AUROC comparison with bootstrap confidence intervals''; Figure~\ref{fig:external_baselines_metrics}, ``Metric profile for Full LCAD--RASA and external baselines under the locked held-out protocol''). Clinical-only baselines used age, HPV, and TCT only; centre identifiers, report text, pathology-like attributes, and labels were excluded from their inputs.

Full LCAD--RASA was competitive with several conventional baselines, including clinical-only logistic regression (AUROC 0.830), colposcopy-only embedding MLP (AUROC 0.827), and late-fusion MLP (AUROC 0.824). It did not, however, dominate all stronger external baselines. Clinical-only HistGradientBoosting reached AUROC 0.862, the cross-attention multimodal transformer reached AUROC 0.847, and the CLIP-style contrastive multimodal baseline reached AUROC 0.888. In the corrected paired bootstrap recheck, the contrastive baseline exceeded Full LCAD--RASA in AUROC by 0.056 (95\% bootstrap interval 0.024--0.092 in favour of the comparator; two-sided $p=0.002$), whereas differences against clinical-only logistic regression, clinical-only HistGradientBoosting, late-fusion MLP, and cross-attention transformer were not statistically conclusive (Table~\ref{tab:paired_bootstrap_recheck}, ``Corrected paired bootstrap AUROC recheck''; Figure~\ref{fig:external_pairwise_recheck}, ``Corrected paired bootstrap AUROC differences between Full LCAD--RASA and external baselines'').

These results sharpen the claim boundary of the study. LCAD--RASA should not be presented as a universally superior pure discriminative classifier over all standard multimodal baselines. Its main contribution is instead report-aware multimodal learning: it converts incomplete report supervision into structured, modality-grounded semantic anchors; provides section-level alignment targets; and supports auditable pseudo-report generation under sparse report supervision. The strong contrastive result also motivates future work combining report-section supervision with contrastive multimodal pretraining. To examine whether the report-aware behaviour was internally coherent, we next analysed component-level ablations.

\subsection{Ablation results linked the observed gains to section-level alignment and operating-point behaviour}
\label{subsec:ablation_results}

The ablation experiments assessed whether the observed performance pattern was consistent with the proposed report-aware alignment mechanism. Removing section alignment substantially reduced direct alignment quality, with macro MRR decreasing from 0.049 for Full LCAD--RASA to 0.022 (Table~\ref{tab:theme1_alignment_ablation}, ``Direct RASA alignment ablation using modality-section retrieval''). In the held-out comparison, the LCAD variant without section alignment retained a similar AUROC of 0.826, but its F1 score decreased to 0.510, compared with 0.611 for Full LCAD--RASA (Table~\ref{tab:main_comparison}, ``Held-out test performance at validation-selected thresholds''). Thus, section-level alignment appears to affect threshold-dependent error balance and operating-point behaviour, even when rank-based discrimination remains similar.

Additional modality, quality-control, alignment-weight, and component ablations are reported in Supplementary Tables~S1 and S3--S5 and Figure~\ref{fig:rasa_lambda} (``Alignment-weight sweep for the RASA objective''), Figure~\ref{fig:rasa_component_boxenplot} (``Modality ablation and RASA component ablation''), and Figure~\ref{fig:qc_ablation} (``LCAD quality-control ablation''). These supplementary ablations were used as diagnostic analyses rather than as final-model superiority tests, because several variants were run under fixed quick-budget settings and some variants produced point estimates close to the full model. The QC ablation in particular did not materially change downstream AUROC/F1 under the current retrospective configuration, so QC is treated here as a procedural safeguard for pseudo-report generation rather than as a major independent driver of diagnostic performance. Together, the ablations support the internal coherence of the RASA design, but they should not be interpreted as causal proof that any single component independently determines the full diagnostic behaviour. Because ablation alone cannot exclude the possibility of fixed-template generation, we further tested whether report sections responded selectively to perturbations of their corresponding modality evidence.

\subsection{Perturbation experiments distinguished modality-grounded decoding from fixed-template generation}
\label{subsec:perturbation_results}

A key validity concern was that generated sections might follow fixed templates despite satisfying section-level metrics. We therefore perturbed modality evidence and measured whether the corresponding report sections degraded selectively. The upgraded perturbation matrix showed that removal or disruption of a modality produced the largest degradation in the matching section (Table~\ref{tab:perturbation_matrix}, ``Upgraded perturbation sensitivity matrix''; Figure~\ref{fig:perturbation_heatmap}, ``Perturbation condition by report-section similarity matrix''; Figure~\ref{fig:perturbation_lineplot}, ``Section-level degradation profiles under modality perturbations''; Figure~\ref{fig:risk_delta}, ``Risk-score shifts under report-generation perturbations''; Figure~\ref{fig:theme1_perturbation}, ``Theme 1 upgraded perturbation sensitivity matrix''). Masking OCT produced a 0.743 drop in the OCT findings section, masking colposcopy produced a 1.000 drop in the colposcopy findings section, and masking clinical instruction produced a 0.833 drop in the clinical-context section.

Label-only inference produced the largest overall report drop of 0.466 and the largest absolute risk shift of 0.371. This result indicates that the model's report behaviour was not fully recoverable from class labels alone. The section-specific degradation pattern therefore provides evidence of modality-grounded decoding and helps distinguish LCAD--RASA from label-copying or schema-filling baselines. The evidence remains perturbational rather than causal; accordingly, we do not claim that the generated reports constitute causal clinical explanations. We then examined whether the observed performance pattern held across reference-report strata and across centres.

\subsection{Reference-stratified and cross-centre analyses exposed the limits of pooled performance}
\label{subsec:loco_results}

Reference-stratified testing examined whether held-out performance was confined to cases with archived reference reports. Full LCAD--RASA achieved AUROC 0.824 among cases with reference reports ($n=108$) and AUROC 0.850 among cases without reference reports ($n=180$). This result suggests that the model's held-out behaviour was not restricted to the subset with archived reports, although the analysis remains retrospective and conditioned on the available cohort composition.

Cross-centre evaluation presented a more restrictive picture. Under strict leave-one-centre-out retraining with the fixed quick-budget protocol, held-out AUROC was 0.702 for Enshi, 0.648 for Jingzhou, and 0.382 for Xiangyang (Supplementary Tables~S2 and S2b; Figure~\ref{fig:loco_heatmap}, ``Leave-one-centre-out evaluation: AUROC heatmap and forest-style centre summary''; Figure~\ref{fig:loco_forest}, ``Strict LOCO centre-wise behaviour under fixed quick-budget retraining''). Eval-only LOCO using a global checkpoint is reported separately and is not pooled with strict LOCO retraining. These results show that centre shift remains a material limitation. The evidence therefore supports report-aware fusion under retrospective held-out testing, but does not support an unrestricted external-generalisation claim across all centres.

\subsection{Robustness, safety, and scalability analyses supported retrospective feasibility}
\label{subsec:robustness_safety_scalability}

Finally, we assessed whether the proposed workflow was reproducible and computationally feasible at cohort scale. Robustness and safety analyses included masking validation, modality ablations, quality-control ablations, multi-seed stability, report-level safety metrics, and runtime summaries (Supplementary Tables~S7--S11; Supplementary Figure~\ref{fig:supplementary_robustness}, ``Supplementary robustness checks: masking validation and multi-seed stability''). The pipeline processed approximately 137k images, and the publishable model checkpoints were approximately 80 MB, indicating that the workflow is feasible for retrospective large-scale multimodal analytics.

Taken together, the experiments show that LCAD--RASA addresses a specific limitation of multimodal cervical screening datasets: report-level semantic supervision is sparse even when image and structured clinical evidence are available at scale. Structured pseudo-reports provided modality-grounded semantic anchors, improved report-scarcity surrogates, and supported numerically stronger held-out performance in the locked retrospective test set. At the same time, strict cross-centre evaluation revealed unresolved centre-shift sensitivity. The resulting interpretation is therefore deliberately bounded: LCAD--RASA provides a retrospective, report-aware fusion framework for exploiting incomplete semantic supervision, but further external validation and prospective assessment are required before clinical deployment claims can be made.
